%%%%%%%%%%%%%%%%%%%%%%%%%%%%%%%%%%%%%%%%%
% Template formal para se utilizar na equipe de propulsão e tecnologia aeroespacial (EPTA)
% LaTeX Template
% Version 1.0 (16/07/2019)
%
% This template was downloaded from:
%
% Original author:
% Peter Wilson (herries.press@earthlink.net) with modifications by:
% Vel (vel@latextemplates.com)
% Vel (feliperibeiro.ufu@gmail.com)
% Vel (mairaf_miranda@hotmail.com)
%
% License:
% CC BY-NC-SA 3.0 (http://creativecommons.org/licenses/by-nc-sa/3.0/)
%
% Conselhos para se lidar com este template:
% - Evitar alterar códigos referentes à formatação, com exceção de se copiar e colar
%          o código para exibição de imagens, tabelas e equações.
% - Procurar sempre manter as imagens utilizadas no documento em uma pasta dedicada,
%          de forma que, ao se referenciar a imagem, se referencie o caminho para 
%          ela, como feito com a "imagem_de_missao".
%
%
%%%%%%%%%%%%%%%%%%%%%%%%%%%%%%%%%%%%%%%%%

%----------------------------------------------------------------------------------------
%	Pacotes utilizados.(evitar mudar, se for necessária alguma mudança, comunicar gerencia)
%----------------------------------------------------------------------------------------

\documentclass[a4paper, 12pt,  oneside,a4paper,english,french,spanish,brazil]{book}     % Para folha A4, default 11pt tamanho
\usepackage[utf8]{inputenc}                                               % Acentos
\usepackage[brazil]{babel}                                                % Biblioteta para português
\usepackage{amssymb}                                                      % Biblioteta para simbolos
\usepackage{graphicx}                                                     % Biblioteta para figuras
\usepackage{tikz}                                                         % Pacote gráfico
\usepackage{amsmath}                                                      % Biblioteta para simbolos
\usepackage[T1]{fontenc}                                                  % encode para acentos 
\usepackage{fouriernc}                                                    % New Century Schoolbook font
\usepackage{xcolor}                                                       % Controle de cores
\usepackage[many]{tcolorbox}                                              % Controle de cores
\usepackage{amsthm}                                                       % Suport matemático ams
\usepackage{physics}                                                      % Suport em notações
\usepackage{gensymb}                                                      % Suport em notação
\usepackage{amsfonts}                                                     % Suport em fontes ams
\usepackage[left=1.00cm, right=1.00cm, top=3cm, bottom=2.50cm]{geometry}  % Parametragem geométrica
\usepackage{caption}                                                      % Legenda para imagens
\usepackage{subcaption}                                                   % Legenda para figuras multiplas
\usepackage{capt-of}                                                      % Suport com legendas
\usepackage{float}                                                        % Suport para maiores floats  
\usepackage[figurename=Fig.]{caption}                                     % Suporte 2 titulos legendas
\usepackage{tabularx}                                                     % suporte com tabelas
\usepackage{ragged2e}                                                     % suporte localização texto
\usepackage{mathrsfs}                                                     % Fontes para equações
\usepackage{makecell}                                                     % Formatação de tabelas
\usepackage{textpos}                                                      % Facilita posicionamento de forma absoluta
\usepackage[hyphens]{url}                                                 % Possibilita uso de hifen em links
\usepackage[makeroom]{cancel}                                             % Desconsidera limites para retas
\usepackage{empheq}                                                       % Possibilita emoldurar fórmulas
\usepackage[colorlinks=true,linkcolor=blue,urlcolor=blue,bookmarksopen=true, bookmarksnumbered, pdfencoding=auto]{hyperref}                                               % links clicaveis em PDF
\usepackage{fancyhdr}                                                     % Ajuda na criação de footers e headers
\usepackage{siunitx}                                                      % General suport
\usepackage{enumitem}                                                     % Controle sobre layout
\usepackage{cancel}                                                       % Riscar coisas
\usepackage{bookmark}                                                     % Hiperlink em PDF
\usepackage{versions}                                                     % Devido às verções
\usepackage[pages=some,scale=1,angle=0,opacity=1]{background}             % Imagem de Background
\usetikzlibrary{decorations.pathmorphing}                                 % Permite utilização direta fonte abaixo
\usetikzlibrary{shapes.callouts,positioning}                              % Balões   
\usetikzlibrary{shapes.geometric,arrows,shadows}                          % flechas
\usepackage{frcursive}                                                    % Fonte cursiva
\usepackage{multirow}                                                     % Suporte tabelas

\setlength\parindent{3ex}                                                 % Parametragem geométrica
\setlength{\parskip}{0.25em}                                              % Parametragem geométrica
\setlength{\fboxsep}{10pt}                                                % Espaço ao redor dos quadros
\definecolor{zelena}{RGB}{0,100,0}                                        % Define cores em estruturas
\tcbuselibrary{theorems}                                                  % Cor nas caixas                           
\tcbuselibrary{breakable}                                                 % Particioina coisas
\pagestyle{fancy}  
\fancyhead[RO , RE]{\footnotesize 
	\begin{minipage}[b]{0.435\linewidth}\flushright\leftmark\end{minipage}}  % cabeçalho direito ímpar=título do capítulo
\fancyhead[LO , LE]{\footnotesize 	
	\begin{minipage}[b]{0.435\linewidth}\flushleft\rightmark\end{minipage}} % cabeçalho esquerdo = nome da seção
\fancyfoot[CE,CO]{\thepage}                                               % Nota de pé
\allowdisplaybreaks                                                       % Possibilita alguns comandos de controle
\setcellgapes{5pt}                                                        % Formtatção de tabelas
\definecolor{eggshell}{rgb}{0.94, 0.92, 0.84}                             % Definição de cor 



\setlength\headheight{27.05003pt}
\renewcommand\headrule{\hrulefill
	\raisebox{10.1pt}[-10pt][-10pt]{\plogo}\hrulefill}                      % Headers
\newcommand{\plogo}{\fbox{  $EPTA$   }}                                   % sigla equipe 
\newcommand\BackImage[2][scale=1]{\BgThispage\backgroundsetup{
		contents={\includegraphics[#1]{#2}}}}                             % cria a função background
\newtcbox{\caixaeq}[1][]{nobeforeafter,math upper,tcbox raise base,enhanced,
	colframe=black!50!black,colback=eggshell!40!white,arc=4pt,	
	boxrule=1pt,drop fuzzy shadow,#1}                                     % Moldura para fórmulas

\DeclareMathOperator{\tg}{tg}                                             % Declara construção mat. 
\DeclareMathOperator{\cotg}{cotg}                                         % Declara construção mat.
%%%%%%%%%%%%%%%%%%%%%%%%%%%%%%%%%%%%%%%%%

\frontmatter
\date{}
%\thispagestyle{empty}
\date{\today}
\sloppy

%----------------------------------------------------------------------------------------
%    INÍCIO DO DOCUMENTO    (a partir daqui pode editar sem grandes complicações)
%----------------------------------------------------------------------------------------
\begin{document} 

%----------------------------------------------------------------------------------------
%	PÁGINA DE TÍTULO
%----------------------------------------------------------------------------------------

 \begin{titlepage} % Suprime cabeçalhos e notas na base da folha.

		\BackImage[scale = 1]{midia/epta_projeto_template}% imagem de background
	
	\begin{minipage}[h!]{0.37\textwidth}
		.%\textcolor{white}{.}% Para ocupar espaço (some times nothing is everything)
	\end{minipage}
	\begin{minipage}[h!]{0.6\textwidth}
	
		\centering % centrliza os textos na página
	
		\scshape % Usa formatação pequena 
	
		%------------------------------------------------
		%	Título
		%------------------------------------------------
	
		\rule{\textwidth}{1.6pt}\vspace*{-\baselineskip}\vspace*{2pt} % Linha horizontal grossa
		\rule{\textwidth}{0.4pt} % Linha horizontal fina
	
		\vspace{0.75\baselineskip} % espaço sobre o título
	
		{\LARGE Desafio Entertainment\\} % Título
	
		\vspace{0.75\baselineskip} % espaço abaixo do título
	
		\rule{\textwidth}{0.4pt}\vspace*{-\baselineskip}\vspace{3.2pt} % Linha horizontal fina
		\rule{\textwidth}{1.6pt} % Linha horizontal grossa
	
		\vspace{2\baselineskip} %espaço em branco abaixo do título
	
		%------------------------------------------------
		%	Subtítulo
		%------------------------------------------------
	
		DESAFIO PROJETADO PARA FUTURA TRANSIÇÃO INTERNA DA EQUIPE DE PROPULSÃO E TECNOLOGIA AEROESPACIAL % descrição adicional (sub-título)
	
		\vspace*{3\baselineskip} % espaço abaixo do subtítulo
	
		%------------------------------------------------
		%	Autores(s)
		%------------------------------------------------
	
		Desenvolvido Por
	
		\vspace{0.5\baselineskip} % Espaço antes de autores
	
		{\scshape\Large Vicenzo De Marco Olivalves\\} % Lista de autores (ordem alfabética)
	
		\vspace{0.5\baselineskip} % Espaço após os autores


            \vspace{0.5\baselineskip}
            
		\textit{Universidade Federal de \\ Uberlândia} % Entidades envolvidas
	
		\vspace{14.8\baselineskip} % Espaço (diminuir se houver mais autores)
	
		{\large Equipe de Propulsão e Tecnologia Aeroespacial} % equipe
		
		\vspace{0.3\baselineskip} % Espaço embaixo da logo
		
		Entertainment  % Núcleo análogo
		
		\vspace{0.3\baselineskip} % Espaço embaixo da logo
		
		\today % Data de compilação
		
	\end{minipage}

\end{titlepage}

%----------------------------------------------------------------------------------------
%	Desenvolvimento...
%----------------------------------------------------------------------------------------
\mainmatter{}
 % Índice
\thispagestyle{empty}    % Sem numeração na página do índice
           % Lista de figuras 
\thispagestyle{empty}     % Sem numeração na página da lista de imagens
%  <------------  Esse comando cria um capítulo

\newpage

\setcounter{page}{1}

\section*{Desafio}

O desafio de transição interna para os futuros integrantes da área do Entertainment da Equipe de Propulsão e Tecnologia Aeroespacial (EPTA) é apresentado neste documento.

Esse formato avaliativo foi moldado de acordo com as competências esperadas dos membros e as necessidades reais da área em questão, almejando sempre a criação estruturada de novos projetos para o ano de 2026.

\section*{Primeira Parte: Expandindo o Bestiário}

Neste quesito, o participante deve idealizar dois novos inimigos temáticos, pertencentes ao Bioma Cristalizado do jogo. 

O primeiro inimigo deve atuar na classe de "Suporte (Debuffer)". Ele não deve possuir uma função de hostilidade direta ou causar dano letal imediato ao jogador, mas sim possuir a habilidade tática de roubar, suprimir ou remover temporariamente um dos \textit{upgrades} atrelados à infusão atual do personagem. 

O segundo inimigo deve ser caracterizado como um \textit{mob} "menor" --- embora não obrigatoriamente seja de pequeno porte visual. Em contrapartida, ele exige mecânicas e padrões agressivos mais simples, atuando de maneira que justifique seu aparecimento orgânico logo nas fases iniciais de exploração.

O objetivo deste desafio é exercitar a apresentação de mecânicas inovadoras que consolidam uma identidade balanceada para o jogo. Pede-se veementemente rigoroso foco às regras, interações e concepções já estabelecidos na base do projeto original, não poupando criatividade narrativa para afastar o gameplay de conceitos que seriam comumente repetitivos. 

Deve ser definida uma identidade visual coesa para ambos os inimigos. São cruciais: a apresentação de sketches (ou descrições textuais ricas e detalhadas) e uma paleta de cores consistente de forma que as criaturas pareçam orgânicas na Lore (a história mundial do jogo).

\section*{Segunda Parte: A Sinergia do ''Artesanato" Macabro}

A principal mecânica de progressão e sobrevivência adotada por nossa estrutura de Game Design subverte o tradicional acúmulo de Experiência (XP). Em nosso contexto, o protagonista fortalece o seu arsenal acoplando pedaços corporais dos inimigos caídos à sua arma central de mineração através da conversão de uma energia orgânica denominada \textit{Essência da Vida}.

Isto posto, a partir dos dois inimigos recém-criados na Primeira Parte, pede-se que o participante desenvolva \textbf{três (3) ''Partes do Corpo'' (Upgrades/Loot)} fornecidas como recompensa (drops) ao se abater essas criaturas, destinadas a se integrarem sem conflitos matemáticos ao sistema:

\begin{itemize}
    \item \textbf{1 Parte de Tier 1 (Natureza Comum):} Orientada a modificações numéricas leves e utilitárias no jogador.
    \item \textbf{1 Parte de Tier 3 (Natureza Rara):} Abordando elevações robustas de atributos numéricos aliadas a algum tipo de comportamento ou \textit{status mod} exclusivo. 
    \item \textbf{1 Parte de Tier 4 (Natureza Lendária):} Que, de maneira invasiva nas regras fundamentais, altera de forma ativa a conduta do jogador e o desenrolar das dinâmicas.
\end{itemize}

Para que seja aprovado com máxima eficácia, ofereça também os seguintes itens narrativos: um batismo adequado condizente com a biologia das criaturas recém-criadas para cada um, origens plausíveis em sua concepção morfológica e, é claro, a especificação das estatísticas exatas modificadas. 

\textbf{Restrição de Game Design:} O item de Tier Lendário (Tier 4) categoricamente é proibido de possuir qualquer tipo de relacionamento com ganhos simples de dano por ataque. Ele precisa influenciar de forma inusitada e sinérgica as mecânicas de movimento evasivo do astronauta (como a esquiva ou um "Dash"), criando novos patamares de gameplay imprevisível, exigindo do candidato criatividade em sua documentação analítica para atestar o quão divertido e inovador isso resultaria na percepção do Operador/Jogador.

\section*{Terceira Parte: Quebrando a Rotina}

Dentre a lista atualizada dos defensores do sistema ecológico base (o Bestiário), há o **Golem** (com o custo gerencial local categorizado como um oponente robusto tipo Tanque). O seu papel primário é ser uma ameaça letal em ataques a curtíssima distância (\textit{Melee}) e de reações propositalmente cadenciadas e lentas.

Em protótipos de combates gerados aleatoriamente em grandes arenas desobstruídas, observou-se uma grave distorção e comodismo comportamental — a falha na dinâmica. O jogador encontrou grande simplicidade e tédio tático ao combater fugindo de costas repetidamente pelas laterais do ambiente e rotacionando a sala aos poucos sem sofrer riscos críveis.

Uma resolução conservadora sugeriria simples mutações no comportamento central, como aumentar passivamente a sua velocidade de locomoção direcional ou anexar ataques remotos por projéteis lineares. Todavia, isso desvirtuaria a natureza intrínseca do oponente estipulada no Game Design inicial.

O participante deve atuar com autoridade resolutiva para recusar essas vertentes conservadoras e implementar soluções inventivas que invalidem esta passividade sistemática de fuga paralela do jogador. Exige-se uma dinâmica que instaure urgência sem ceder mobilidade anormal ao nosso adversário massivo e pesado — preservando uma experiência assustadora. Pede-se ampla descrição e a justificativa técnica das dinâmicas propostas deste obstáculo para o gameplay geral da ''Arena''.

\section*{Bônus: Identidade Nominal}

O projeto atualmente possui um alicerce narrativo estabelecido: trata-se de um universo mecânico e biológico hostil centrado na sobrevivência de um astronauta e o seu cão cibernético (Eptinho) em um exoplaneta alienígena. Contudo, a obra ainda carece de uma de suas características fundamentais: um Título Oficial.

Como diretiva complementar, pede-se que o participante idealize sugestões para a nomenclatura formal do jogo. Procura-se um batismo que encapsule e transmita ao público alvo as temáticas de sobrevivência de forma sintética, profunda e memorável. Será muito valorizada a justificativa do nome escolhido bem como seu impacto interpretativo no mercado focal.

\newpage

\setcounter{page}{2}

\begin{center}
\fbox{EPTA}
\end{center}

\vspace{0.5cm}

\section*{Observação}

Tem-se conhecimento de que esse é um desafio de cunho profundo, constituído por um alto volume de requisitos mentais a serem geridos e resolvidos em diferentes âmbitos na engenharia dos sistemas de divertimentos e que isso toma considerável tempo; posto isto, compreendem-se perfeitamente eventuais lacunas no escopo submetido. A análise principal do membro dar-se-á justamente pela perspicácia nas abordagens fragmentadas em sua composição macro; portanto, focar intensivamente nas qualidades intelectuais e na profundidade dos trechos em que houver investimento deliberado resultará em uma ótima avaliação qualitativa.

Ressalta-se fundamentalmente o pressuposto norteador do teste focado inteiramente na idealização imaginativa da jogabilidade e do ecossistema estrutural; desta ótica o perfil humano destacado como preferencial a nossa vertente EPTA Entertainment são as visões singulares pautadas na inconfundível capacidade técnica lúcida e criacional.

Qualquer dúvida cabível em relação aos preceitos listados, não se hesite em firmar diálogos proativos requisitando suportes operacionais ou brainstormings, reportando diretamente ao gerente atual da área:

Vicenzo De Marco Olivalves

(34)991857053

\end{document}
